\section{A tie with no bond drift}
\label{sec:h-rewire}

A degree-preserving re-pairing changes the local isomorphism type at
\[
\Delta E=-2(B_{\mathrm{add}}-B_{\mathrm{remove}}).
\]
When the two completions of a \(4\)-cycle have the same \(\Delta E\), the bond
does not select.
REWIRE0 finds the same tie, or an absence, under the frozen exact principles,
persistently in the linear size.
A branching ratio would be a measure on that tie.
No such measure is forced by the earned constraints.
You could still compare the two rewired graphs by the single number
\(\Delta E\), on a substrate other than \(J_2\).
A splitting that stays finite as the linear size grows would be a scale in
the same units as \(B\).
A splitting that closes leaves the tie flat.
That number is the same \(B\) as in Section~\ref{sec:bond}.
The current of either edge does not enter \(\Delta E\).
Moving along the circle of Proposition~\ref{prop:circle} is a different
operation: it changes the endpoint amplitudes, and
Proposition~\ref{prop:circle-energy} says the contraction energy then
moves unless the exclusive-neighbour imbalance \(V\) vanishes.
